\documentclass{beamer}
\usepackage{amsmath, amsfonts}
\usepackage{listings}
 
 
\usepackage{amsmath, amsfonts}
\usepackage{listings}
\usepackage{xcolor}
\usepackage{graphicx}
\usepackage{booktabs}
\usepackage{tikz}
\usepackage{hyperref}
\usepackage{multicol}
\usepackage{lmodern}  
% Polish language and UTF-8 support
\usepackage[T1]{fontenc}
\usepackage[utf8]{inputenc}
\usepackage[polish]{babel}
\usepackage{caption}
\usepackage{libertine}
\usepackage[libertine]{newtxmath}
\usepackage{algorithm}
\usepackage{algpseudocode}
 
\newcommand{\EFL}{Erdős-Faber-Lovász }
 
\title[Enhanced Beamer Demo]{\Large{DOŚWIADCZENIE ISTNIENIA} \\ \vspace{1em}
Czy autentyczność egzystencji jest problemem ontycznym czy ontologicznym?}
\author{Jakub Wingorowicz, Antoni Wojtasik, Szymon Wojtulewicz}
\date{15 czerwca 2026}
 
\usetheme{default}
 
\usefonttheme{professionalfonts}
 
 
\definecolor{accent}{RGB}{130,160,255}
\definecolor{blockbg}{RGB}{240,240,245}
\definecolor{blockborder}{RGB}{220,220,220}
 
\setbeamertemplate{navigation symbols}{}
\setbeamertemplate{blocks}[default]
 
 
\setbeamercolor{block title}{fg=black,bg=blockbg}
\setbeamercolor{block body}{bg=blockbg}
% \setbeamerfont{block title}{series=\sffamily\bfseries}
 
\definecolor{codebg}{RGB}{250,250,250}
\definecolor{keywordcolor}{RGB}{70,110,200}
\definecolor{stringcolor}{RGB}{200,130,130}
\definecolor{commentcolor}{RGB}{140,180,140}
 
\lstset{
  backgroundcolor=\color{codebg},
  basicstyle=\ttfamily\footnotesize,
  keywordstyle=\color{keywordcolor}\bfseries,
  stringstyle=\color{stringcolor},
  commentstyle=\color{commentcolor}\itshape,
  frame=single,
  breaklines=true,
  showstringspaces=false,
  tabsize=2
}
 
 
\begin{document}
 
\frame{\titlepage}
 
% ---------------------------------------------------------------
% CZĘŚĆ I — RAMA POJĘCIOWA
% ---------------------------------------------------------------
 
\begin{frame}{Teza}
  \large
  \begin{block}{}
    Filozoficzne wyjaśnienie tego, \emph{jak rozpoznać prawdziwe ja},
    jest niemal niemożliwe do sformułowania.
  \end{block}
  \vspace{1em}
  Każda próba wskazania prawdziwego ja prowadzi do jednego z trzech wyników:
  \begin{itemize}
    \item regresu,
    \item błędnego koła,
    \item kontrowersyjnych założeń metafizycznych.
  \end{itemize}
\end{frame}
 
\begin{frame}{Trzy podejścia do autentyczności}
  \huge {
    \begin{enumerate}
      \item Ontologiczne
      \item Ontyczne
      \item Etyczne
    \end{enumerate}
  }
\end{frame}
 
\begin{frame}{Krytyka podejścia ontologicznego}
  \huge {
    \begin{enumerate}
      \item Problem Własnościowy
      \item Problem Logiczny
      \item Problem Czasowy
    \end{enumerate}
  }
\end{frame}
 
% ---------------------------------------------------------------
% TRZY PROBLEMY — ROZWINIĘCIE
% ---------------------------------------------------------------
 
\begin{frame}{Problem Czasowy --- autentyczność jako oryginalność historyczna}
  \begin{itemize}
    \item Coś jest autentyczne, gdy zachowuje ciągłość z bytem historycznym.
    \item Kategoria temporalna: sens pojawia się dopiero z upływem czasu
          i przy możliwości fałszerstwa.
    \item Przeniesiona na osobę rodzi pytanie: \emph{który moment życia} jest
          punktem odniesienia?
  \end{itemize}
  \vspace{0.5em}
  \begin{block}{Dylemat}
    Bardziej autentyczny jest człowiek młody i spontaniczny czy dojrzały,
    po procesie samorozwoju? Autentyczność nie jest prostym powrotem do
    przeszłości --- a tę intuicję definicja podważa.
  \end{block}
\end{frame}
 
\begin{frame}{Problem Własnościowy --- autentyczność jako przeciwieństwo fałszu}
  \begin{itemize}
    \item Prawdziwa perła vs.\ podróbka: decyduje źródło i natura, nie wiek.
    \item Zastosowane do osoby zakłada istnienie \emph{fałszywego ja},
          podszywającego się pod właściwą osobę.
    \item Sytuacja możliwa tylko w eksperymentach myślowych: klony,
          bliźniacze światy, doskonałe kopie.
  \end{itemize}
  \vspace{0.5em}
  \begin{block}{Błędne koło}
    Aby wskazać instancję autentyczną, potrzebujemy cechy odróżniającej ---
    zwykle pierwszeństwa czasowego lub genealogicznego. To powrót do
    \emph{problemu czasowego}.
  \end{block}
\end{frame}
 
\begin{frame}{Problem Logiczny --- autentyczność jako wierność sobie}
  \begin{itemize}
    \item Najpowszechniejsze znaczenie: być uczciwym i wiernym sobie.
    \item Paradoks (Berman): jeśli autentyczność to \emph{bycie sobą}, to czym
          jest nieautentyczność? W jaki sposób można nie być sobą?
    \item Idea narusza zasadę \(A = A\) --- każdy jest sobą niezależnie od tego,
          czy jest szczery, czy zakłamany.
  \end{itemize}
  \vspace{0.5em}
  \begin{block}{Różnica wobec rzeczy}
    Dla pereł i zabytków wiemy, jakie dane rozstrzygają o autentyczności.
    Dla osób takich danych nie ma --- ocena staje się aktem interpretacji
    (problem epistemiczny \emph{i} metafizyczny).
  \end{block}
\end{frame}
 
\begin{frame}{Problem zmiennego i stającego się ja}
  \begin{itemize}
    \item Rousseau i Nietzsche: człowiek nie jest bytem statycznym.
    \item W przeciwieństwie do przedmiotów osoba \emph{nieustannie staje się} sobą.
    \item Autentyczność zdaje się wymagać trwałej tożsamości, której można
          być wiernym\ldots
    \item \ldots tymczasem egzystencja to ciągła zmiana i reinterpretacja
          własnej przeszłości.
  \end{itemize}
  \vspace{0.5em}
  Im mocniej podkreślamy dynamikę życia, tym trudniej wskazać niezmienne
  centrum osobowości pełniące funkcję ,,prawdziwego ja''.
\end{frame}
 
% ---------------------------------------------------------------
% EGZYSTENCJALIZM
% ---------------------------------------------------------------
 
\begin{frame}{Odpowiedź egzystencjalizmu (Heidegger, Sartre)}
  \begin{itemize}
    \item Odrzucenie tradycyjnej idei ukrytej esencji --- przy zachowaniu
          pojęcia autentyczności.
    \item Autentyczność to nie odkrycie prawdziwego ja, lecz \emph{sposób życia}.
    \item Człowiek autentyczny uznaje własną wolność, odpowiedzialność
          i skończoność.
    \item Heidegger: chodzi o wierność strukturze ludzkiego bycia (Dasein),
          a nie indywidualnej naturze.
  \end{itemize}
\end{frame}
 
\begin{frame}{Granice egzystencjalizmu}
  \begin{itemize}
    \item Guignon: zachowany zostaje motyw \emph{ukrytego źródła}, któremu należy
          pozostać wiernym.
    \item Źródłem nie jest jednostka, lecz czasowa struktura istnienia jako takiego.
    \item Efekt: esencjalizm przeniesiony z poziomu jednostki na poziom gatunku.
  \end{itemize}
  \vspace{0.5em}
  \begin{block}{Paradoks}
    Egzystencjalizm jednocześnie odrzuca ideę prawdziwego ja
    i potrzebuje jej odpowiednika, by zachować sens autentyczności.
  \end{block}
\end{frame}
 
\begin{frame}{Autentyczność a wartości obiektywne}
  \begin{itemize}
    \item Taylor, Varga, Cottingham: autentyczność jako ideał charakteru
          zakorzeniony w obiektywnych wartościach.
    \item Pozwala uniknąć radykalnego subiektywizmu --- rozwijanie siebie
          wobec wartości, a nie ekspresja dowolnych pragnień.
  \end{itemize}
  \vspace{0.5em}
  \begin{block}{Pytanie, które pozostaje}
    Dlaczego \emph{ten} zestaw wartości miałby ujawniać moje prawdziwe ja
    bardziej niż jakikolwiek inny? Problem identyfikacji nie znika ---
    przybiera nową formę.
  \end{block}
\end{frame}
 
% ---------------------------------------------------------------
% CZĘŚĆ II — PROPOZYCJA NIEESENCJALISTYCZNA
% ---------------------------------------------------------------
 
\begin{frame}
  \centering
  \vfill
  {\Huge W stronę nieesencjalistycznej\\[0.3em] koncepcji autentyczności}
  \vfill
  {\large Część II}
  \vfill
\end{frame}
 
\begin{frame}{Buddyjska krytyka ,,prawdziwego ja'' --- \emph{anattā}}
  \begin{itemize}
    \item Buddyzm odrzuca ideę wiecznej i niezmiennej jaźni.
    \item ,,Ja'' to nie substancja ani dusza, lecz konwencjonalna nazwa
          dynamicznego zespołu procesów psychofizycznych.
  \end{itemize}
  \vspace{0.5em}
  Cztery cechy stanowiska:
  \begin{enumerate}
    \item nondualizm ontologiczny,
    \item pragmatyzm epistemologiczny (rola praktyki kontemplacyjnej),
    \item tożsamość procesualna, nietrwała, współzależna,
    \item cel: wyzwolenie od cierpienia, nie odkrycie metafizycznej prawdy.
  \end{enumerate}
  \vspace{0.3em}
  Autentyczność zaczyna się tam, gdzie porzucamy iluzję trwałego ja.
\end{frame}
 
\begin{frame}{Pięć skupisk (skandh)}
  Osoba jako zbiór pięciu procesów --- żaden nie jest trwałym rdzeniem:
  \begin{enumerate}
    \item forma cielesna,
    \item uczucia,
    \item percepcja,
    \item formacje mentalne,
    \item świadomość.
  \end{enumerate}
  \vspace{0.5em}
  \begin{block}{Wniosek}
    Jaźń jest \emph{procesem}, a nie rzeczą.
  \end{block}
\end{frame}
 
\begin{frame}{Kognitywistyczny obraz jaźni}
  \begin{itemize}
    \item Brak pojedynczego, centralnego podmiotu sterującego zachowaniem.
    \item Rozproszone, dynamiczne procesy neuronalne integrują informacje
          sensoryczne, motoryczne, emocjonalne i poznawcze.
    \item W mózgu nie ma ,,siedliska'' trwałej tożsamości --- jest ona
          \emph{funkcjonalnym osiągnięciem} systemu poznawczego.
    \item Jaźń narracyjna: opowieść o sobie tworzona dzięki pamięci
          autobiograficznej i językowi.
    \item Sieć trybu domyślnego (DMN): autorefleksja, wspomnienia,
          wyobrażanie przyszłości.
  \end{itemize}
  \vspace{0.3em}
  \footnotesize Korelat neuronalny \(\neq\) ontologiczny fundament.
\end{frame}
 
\begin{frame}{Relacyjny charakter samorozumienia}
  \begin{itemize}
    \item Teoria umysłu: rozumienie siebie i rozumienie innych opiera się
          częściowo na tych samych procesach.
    \item Samowiedza nie jest izolowaną introspekcją --- powstaje w relacjach,
          dialogu, języku i kulturze.
    \item Jaźń jako \emph{model} tworzony przez mózg w celu organizowania
          doświadczeń i koordynowania zachowania.
    \item Tożsamość współkonstruowana przez jednostkę i jej środowisko społeczne.
  \end{itemize}
\end{frame}
 
\begin{frame}{Lekcja postmodernizmu}
  \begin{itemize}
    \item Derrida: tożsamość konstytuuje się przez różnicę i odraczanie znaczeń.
    \item Foucault: podmiot kształtowany przez praktyki dyskursywne
          i relacje władzy.
    \item Butler: tożsamość jako powtarzane akty performatywne reprodukujące normy.
  \end{itemize}
  \vspace{0.5em}
  Bliskość buddyjskiej krytyki substancjalnego ja, ale\ldots
  \begin{block}{Trudność}
    Jeśli podmiot jest wyłącznie efektem dyskursu i władzy, autentyczność
    grozi staniem się pojęciem pustym lub jedynie grą społeczną.
  \end{block}
\end{frame}
 
% ---------------------------------------------------------------
% PROPOZYCJA: AUTENTYCZNOŚĆ JAKO PRAKTYKA (PODEJŚCIE ETYCZNE)
% ---------------------------------------------------------------
 
\begin{frame}{Autentyczność jako praktyka (podejście etyczne)}
  \begin{itemize}
    \item Nie ma pojedynczego, niezmiennego ja do odkrycia.
    \item Ale całkowita rezygnacja z podmiotowości = utrata sensu
          odpowiedzialności, integralności i rozwoju.
  \end{itemize}
  \vspace{0.5em}
  \begin{block}{Propozycja}
    Autentyczność to nie odkrywanie tego, kim jesteśmy, lecz sposób
    \emph{odnoszenia się} do dynamicznego procesu, którym jesteśmy.
  \end{block}
  \vspace{0.3em}
  Obejmuje: krytyczną autorefleksję, otwartość na korektę przekonań,
  świadomość uwarunkowań oraz działanie zgodne z wartościami.
\end{frame}
 
\begin{frame}{Ontyczny czy ontologiczny problem?}
  \begin{itemize}
    \item \textbf{Podejście ontologiczne} --- prawdziwe ja jako byt do odkrycia
          (esencja jednostki lub struktura Dasein). \\
          \footnotesize Upada na problemach: własnościowym, logicznym, czasowym.
    \item \normalsize \textbf{Wspólne założenie} obu wersji: ,,ukryte źródło'',
          któremu trzeba być wiernym.
  \end{itemize}
  \vspace{0.5em}
  \begin{block}{Odpowiedź}
    Autentyczność egzystencji nie jest ostatecznie ani faktem ontycznym,
    ani strukturą ontologiczną --- jest \emph{praktyką etyczną}
    świadomego uczestnictwa w procesie stawania się.
  \end{block}
\end{frame}
 
\begin{frame}{Zakończenie}
  \begin{itemize}
    \item Tożsamość jest \emph{procesem}, a nie substancją.
    \item Zachowujemy to, co cenne w ideale autentyczności: dążenie do
          integralności, szczerości i odpowiedzialności.
    \item Odrzucamy problematyczne założenie metafizycznego ,,prawdziwego ja''.
  \end{itemize}
  \vspace{0.5em}
  \begin{block}{}
    Autentyczności się nie odkrywa --- autentyczność się \emph{praktykuje}.
    Bez metafizycznego esencjalizmu i bez normatywnego nihilizmu.
  \end{block}
\end{frame}
 
\end{document}
 
